\documentclass{ieeeaccess}
\usepackage{cite}
\usepackage{amsmath,amssymb,amsfonts}
\usepackage{algorithmic}
\usepackage{graphicx}
\usepackage{textcomp}
\usepackage{booktabs}
\usepackage{subcaption}
\usepackage{array}
\usepackage{url}
\urlstyle{same}
\def\UrlBreaks{\do\/\do\-\do\_\do\.}
\def\BibTeX{{\rm B\kern-.05em{\sc i\kern-.025em b}\kern-.08em
    T\kern-.1667em\lower.7ex\hbox{E}\kern-.125emX}}

\newcommand{\crackAcc}{0.970}
\newcommand{\crackFone}{0.970}
\newcommand{\crackPrec}{0.969}
\newcommand{\crackRec}{0.971}
\newcommand{\crackImages}{1600}
\newcommand{\crackAreaPx}{150}
\newcommand{\crackCovNeg}{44.9}
\newcommand{\crackCovPos}{9.2}
\newcommand{\crackCompNeg}{7}
\newcommand{\crackCompPos}{1}
\newcommand{\crackSkelNeg}{1728}
\newcommand{\crackSkelPos}{280}
\newcommand{\crackWidth}{16.5}
\newcommand{\crackTreeCV}{0.970}
\newcommand{\crackTopFeat}{skeleton per diagonal}
\newcommand{\crackTopImportance}{0.94}
\newcommand{\crackLowImportance}{0.04}
\newcommand{\crackLogCV}{0.981}
\newcommand{\crackCannyFone}{0.945}
\newcommand{\crackSobelFone}{0.971}
\newcommand{\wbcAcc}{0.596}
\newcommand{\wbcFone}{0.584}
\newcommand{\wbcImages}{1600}
\newcommand{\wbcBest}{LDA}
\newcommand{\wbcLDA}{0.596}
\newcommand{\wbcNB}{0.506}
\newcommand{\wbcNucleusShare}{0.417}
\newcommand{\riceAcc}{0.961}
\newcommand{\riceFone}{0.961}
\newcommand{\riceImages}{1000}
\newcommand{\riceBest}{LinearSVC}
\newcommand{\riceDropped}{0}
\newcommand{\riceFeatures}{24}
\newcommand{\riceMorphOne}{0.961}
\newcommand{\riceHuOne}{0.839}
\newcommand{\riceChainOne}{0.523}
\newcommand{\riceAllOne}{0.955}


\begin{document}
\history{Manuscript received November 13, 2026; revised November 13, 2026;
accepted November 13, 2026. Date of publication November 13, 2026; date of
current version November 13, 2026.}
\doi{}

\title{Digital Image Processing for Surface Defect, Blood Cell, and Grain
Quality Analysis: A Feature-Centered Study of Morphological, Edge, and
Boundary Descriptors}

\author{\uppercase{Muhammad Gauza Faliha}\authorrefmark{1}}
\address[1]{Department of Computer Science and Electronics, Faculty of
Mathematics and Natural Sciences, Universitas Gadjah Mada, Yogyakarta,
Indonesia (e-mail: muhammadgauzafaliha@mail.ugm.ac.id)}

\tfootnote{This work was prepared as the fourth assignment of the Digital Image
Processing course, Department of Computer Science and Electronics, Faculty of
Mathematics and Natural Sciences, Universitas Gadjah Mada, supervised by
Wahyono, S.Kom., Ph.D. No external funding was received.}

\markboth{Muhammad Gauza Faliha: Digital Image Processing for Surface Defect,
Blood Cell, and Grain Quality Analysis}{Muhammad Gauza Faliha: Digital Image
Processing for Surface Defect, Blood Cell, and Grain Quality Analysis}

\corresp{Corresponding author: Wahyono, S.Kom., Ph.D.}

\begin{abstract}
A recurring question in applied image analysis is not whether a classifier can be
trained, but which description of the image carries the decision. This work
answers that question for three inspection settings solvable without learned
representations: a surface defect, a cytological typing task, and an agricultural
grading task. For each we survey the morphological, edge, and representation
operation families, implement a complete pipeline, and ablate the descriptor set.
Concrete cracks are recovered by gradient magnitude and Otsu thresholding, cleaned
by opening and closing, filtered by component area, and described by ten shape
features of which logged skeleton length and component count are decisive; on
\crackImages{} images the pipeline reaches accuracy \crackAcc{} and crack-class F1
\crackFone{}, with a measured width of \crackWidth{} pixels. A foreground-area rule
fails here: median mask coverage is \crackCovNeg{} percent for uncracked against
\crackCovPos{} percent for cracked surfaces, the opposite of the intuitive
direction. White blood cell nuclei are isolated by a guarded ratio-color
enhancement with an ellipse-expanded cell body, and thirteen region and intensity
features with a linear discriminant reach accuracy \wbcAcc{} over four subtypes.
Rice grains are described by nine shape measurements, seven Hu moments, and an
eight-bin chain-code histogram, where shape alone reaches macro-F1 \riceMorphOne{}
while the full \riceFeatures-feature vector falls to \riceAllOne{}. Feature
suitability is thus set by the task type rather than the object domain, and a
descriptor invariant to the wrong transformation adds variance instead of
information.
\end{abstract}

\begin{keywords}
boundary descriptors, concrete crack detection, edge detection, feature
extraction, Hu moments, mathematical morphology, rice variety classification,
white blood cell segmentation
\end{keywords}

\titlepgskip=-15pt

\maketitle

\section{Introduction}
\label{sec:intro}
\PARstart{D}{igital} image processing (DIP) is often presented as a set of
operators to be memorized, but in practice the operators are chosen by the
question being asked. A pipeline that asks \emph{where} an intensity
discontinuity lies needs a derivative estimator and a decision rule on its
magnitude. A pipeline that asks \emph{which object} is present needs a region
description that is invariant to the object's position. A pipeline that asks
\emph{which category} the object belongs to needs a description whose variability
within a category is smaller than its variability between categories. The three
questions are different, and the third assignment block of the course makes the
distinction concrete by separating morphological processing, edge-based feature
extraction, corner and line extraction, and edge-based feature representation.

This paper takes that separation seriously and tests it. We select three cases
that are solvable by classical techniques alone, one for each of the three
question types above, and for each case we (i) survey the operation families
that the literature uses, (ii) build a complete pipeline from those families,
(iii) measure the pipeline's behaviour on a real public dataset, and (iv)
ablate the feature representation to determine which descriptors actually carry
the decision. The cases are deliberately drawn from unrelated domains, because
the claim under test is that the \emph{task type}, not the \emph{object domain},
determines the suitable feature family.

The first case is surface defect inspection: detecting and quantifying cracks in
concrete. Crack detection has a long classical literature, and critical reviews
of it organize the field around thresholding, edge detection, and morphology
before the recent shift toward learned detectors \cite{mohan2018}. We follow the
thresholding-plus-filtering lineage represented by Talab \emph{et al.}
\cite{talab2016} and measure how much each pipeline stage contributes.

The second case is cytology: segmenting white blood cell nuclei and separating
four subtypes. The reference method of Prinyakupt and Pluempitiwiriyawej
\cite{prinyakupt2015} is a compact classical pipeline, a color-ratio enhancement
for the nucleus, ellipse fitting for the cell boundary, and a feature vector of
morphological and intensity statistics classified by linear discriminant
analysis and naive Bayes; it therefore isolates exactly the feature families
this assignment is about.

The third case is agricultural grading: classifying rice varieties from single
grain images. Cinar and Koklu \cite{cinar2019} describe the morphology-and-moment
feature set that we adopt, and a later deep-learning study on the same data
\cite{koklu2021} provides the contrast that motivates keeping the classical
branch in scope here.

The contributions of the paper are:
\begin{itemize}
\item a complete, reproducible classical pipeline for crack detection with an
explicit stage-by-stage ablation of thresholding, morphological cleanup, and area
filtering, plus a skeleton-based measurement of crack width (Section
\ref{sec:case1});
\item a nucleus segmentation and four-class subtype classification built on a
guarded color-ratio enhancement, with an explicit statement of where our
implementation deviates from the reference method and why (Section
\ref{sec:case2});
\item a three-way feature ablation, morphology versus Hu moments versus chain
code, on single-grain variety classification, which answers the assignment's
question of which descriptors are suitable for a shape-defined class (Section
\ref{sec:case3});
\item a cross-case feature map that ties task type to winning descriptor family
and to the failure mode observed in our own runs (Section \ref{sec:discussion}).
\end{itemize}

Section \ref{sec:method} surveys the operations and descriptors used, Section
\ref{sec:results} reports the three experiments, Section \ref{sec:discussion}
compares them, and Section \ref{sec:conclusion} concludes.

\section{Materials and Methods}
\label{sec:method}

\subsection{Datasets}
\label{sec:datasets}
All three experiments use public datasets, subsampled with a fixed seed so that
every table in this paper can be regenerated. Table \ref{tab:datasets} lists
them. The crack images are the fixed-size two-class collection derived from the
concrete crack imagery of \"Ozgenel \cite{ozgenel2019}, distributed as
$227\times227$ crops; the alternative SDNET2018 collection \cite{dorafshan2018}
is larger but has variable resolution and annotated sub-images, which would add a
resampling variable to a study of morphological operators, so the fixed-size
collection is preferred here. The blood cell images come from the Kaggle blood
cell collection \cite{mooney2018}, which provides four stained subtypes at
$240\times320$. The rice grain images are the single-grain collection of Koklu
\emph{et al.} \cite{koklu2021}, five varieties on a uniform dark background.
Dataset licence and size metadata in Table \ref{tab:datasets} were read from the
provider's programmatic record at the access date. The three collections are
publicly available at
\url{https://www.kaggle.com/datasets/arnavr10880/concrete-crack-images-for-classification},
\url{https://www.kaggle.com/datasets/paultimothymooney/blood-cells}, and
\url{https://www.kaggle.com/datasets/muratkokludataset/rice-image-dataset}
respectively, and the experiment code that consumes them is
\url{https://github.com/} \emph{(repository link to be added on submission)}.

\begin{table}[!t]
\centering
\caption{Datasets used. Counts are the stratified subsample processed; the
parenthesised figure is the full collection size. Public URLs are given in
Section \ref{sec:datasets}.}
\label{tab:datasets}
\footnotesize
\begin{tabular*}{\columnwidth}{@{\extracolsep{\fill}}p{1.1cm}p{1.9cm}c p{2.2cm}@{}}
\toprule
Case & Source & Used & Resolution / licence \\
\midrule
Crack & \"Ozgenel \cite{ozgenel2019} & \crackImages{} (40\,000) &
$227^2$, CC BY 4.0 \\
WBC & Mooney \cite{mooney2018} & \wbcImages{} (12\,500) &
$240\times320$, research use \\
Rice & Koklu \emph{et al.} \cite{koklu2021} & \riceImages{} (75\,000) &
$250^2$, CC0 \\
\bottomrule
\end{tabular*}
\end{table}

\subsection{Common Preprocessing}
Every pipeline begins with colour conversion to a single intensity channel,
because the operators used downstream are defined on scalar images. Conversion
to luminance is followed by Gaussian smoothing, which suppresses sensor noise
before any derivative is taken; the kernel is $5\times5$ with
$\sigma$ derived from the kernel size, and the choice is held constant across
cases so that differences between cases are attributable to the descriptors and
not to the smoothing.

Global thresholding uses Otsu's method \cite{otsu1979}, which chooses the
threshold that maximizes the between-class variance of the intensity histogram.
Its appeal here is that it is parameter-free, so a pipeline built on it has no
tuned threshold to leak information from the test data. Where a case requires
contrast normalization, histogram equalization is applied after the colour-ratio
enhancement rather than to the raw channels, so that the enhancement's dynamic
range is what gets equalized.

\subsection{Morphological Processing}
\label{sec:morph}
Mathematical morphology treats a binary image as a set $A\subseteq\mathbb{Z}^2$
and a structuring element as a set $B$, following the course development of the
operators \cite{handouts}. Translation of $A$ by $z$ is
$A_z=\{a+z\mid a\in A\}$, and the reflection is $\hat{B}=\{-b\mid b\in B\}$. The
two primitives are dilation and erosion,
\begin{equation}
A\oplus B=\{z\mid(\hat{B})_z\cap A\neq\varnothing\},\qquad
A\ominus B=\{z\mid B_z\subseteq A\},
\label{eq:dilero}
\end{equation}
which respectively expand and contract the foreground, the first filling small
gaps and the second removing small protrusions. They are not inverses, and the
two compositions that survive as useful filters are opening and closing,
\begin{equation}
A\circ B=(A\ominus B)\oplus B,\qquad A\bullet B=(A\oplus B)\ominus B .
\label{eq:openclose}
\end{equation}
Opening removes foreground detail smaller than $B$ and is therefore a
speckle-killer; closing removes background detail smaller than $B$ and is
therefore a gap-bridger. Because both are idempotent, applying them twice
changes nothing, which makes their effect on a reported mask auditable.

Two further operators appear in the survey and in our measurements. The
hit-or-miss transform detects a specific local configuration by eroding with one
structuring element and the complement with another,
\begin{equation}
A\circledast (B_1,B_2)=(A\ominus B_1)\cap(A^c\ominus B_2),
\label{eq:hitmiss}
\end{equation}
and thinning is the iterative removal of the configurations that hit-or-miss
identifies as redundant,
\begin{equation}
A\otimes B=A-(A\circledast B),
\label{eq:thin}
\end{equation}
repeated until no change occurs. A skeleton is the limit of thinning with a
sequence of structuring elements; it is the representation we use to convert a
crack mask of arbitrary width into a centreline whose length is measurable.
Equation \eqref{eq:thin} is the basis of the Zhang--Suen procedure
\cite{gonzalez2018} implemented in this work, since the standard Python imaging
stack available here does not expose a thinning routine.

The structuring elements used are ellipses, which are less directional than
squares and therefore do not impose a preferred orientation on anisotropic
objects. The crack case uses a $3\times3$ element for opening and a
$5\times5$ element for closing; the cell case uses $3\times3$, $7\times7$, and
$15\times15$ elements at successive stages; the rice case uses a $7\times7$
element only for hole filling, since a single convex grain needs little
cleanup.

\subsection{Edge Detection}
\label{sec:edges}
An edge is a local maximum of the first derivative or a zero crossing of the
second. The discrete approximations differ in the neighbourhood they use and in
their noise sensitivity. The Roberts operators are $2\times2$ crosses,
\begin{equation}
G_x=\begin{bmatrix}1&0\\0&-1\end{bmatrix},\qquad
G_y=\begin{bmatrix}0&1\\-1&0\end{bmatrix},
\label{eq:roberts}
\end{equation}
and are the cheapest and noisiest. The Sobel operators are $3\times3$ and
include a smoothing direction orthogonal to the differentiation direction,
\begin{equation}
G_x=\begin{bmatrix}-1&0&1\\-2&0&2\\-1&0&1\end{bmatrix},\qquad
G_y=\begin{bmatrix}-1&-2&-1\\0&0&0\\1&2&1\end{bmatrix},
\label{eq:sobel}
\end{equation}
so their gradient magnitude $\lVert\nabla f\rVert=\sqrt{G_x^2+G_y^2}$ is already
partly denoised. The Laplacian-of-Gaussian approach smooths first and then
detects zero crossings of the second derivative, which localizes well but
responds to both polarities of the step.

Canny's formulation \cite{canny1986} makes the three stages explicit: optimal
smoothing, non-maximum suppression along the gradient direction, and hysteresis
thresholding with two levels. Its stated optimality criteria are detection,
localization, and a single response per edge, and those criteria are what the
stage model operationalizes. We therefore report Canny as a comparator for the
crack case but build the measured pipeline on Sobel magnitude, because the
hysteresis thresholds are two free parameters while the stage model in
Section \ref{sec:case1} needs the raw gradient available for thresholding.

The same survey covers corner and line features, which complete the block but
are not used by any case in this paper. Harris and Stephens \cite{harris1988}
characterize a corner by the eigenvalues of the summed second-moment matrix over
a window, and Duda and Hart \cite{duda1972} accumulate votes in a parameter
space so that collinear boundary points reinforce one another. Both are the right
tool when the object of interest is a junction or a straight boundary; none of
our three cases has that structure, so we cite them as surveyed alternatives
rather than as components.

\subsection{Feature Representations}
\label{sec:features}
The assignment's central question is which representation to use, so the
representations are treated as the object of study rather than as an
implementation detail.

\emph{Region descriptors.} For a region of area $A$ and perimeter $P$, the
compactness or roundness
\begin{equation}
R=\frac{4\pi A}{P^2}
\label{eq:roundness}
\end{equation}
equals one for a disc and falls toward zero as the boundary becomes convoluted.
Solidity, the ratio of $A$ to the area of the convex hull, measures concavity;
extent, the ratio of $A$ to the bounding-box area, measures how fully the region
fills its box. Eccentricity follows from the second-order central moments,
\begin{equation}
e=\sqrt{1-\frac{(\mu_{20}-\mu_{02})^2+4\mu_{11}^2}{(\mu_{20}+\mu_{02})^2}},
\label{eq:ecc}
\end{equation}
and separates elongated from circular regions. These five numbers plus area and
perimeter form the morphological part of the crack and grain feature vectors.

\emph{Moments.} Hu's seven moments are normalized central moments chosen to be
invariant to translation, scale, and rotation. Because their dynamic range spans
orders of magnitude, they are used in the log-transformed form
\begin{equation}
h_i=-\operatorname{sgn}(\phi_i)\log_{10}\lvert\phi_i\rvert,
\label{eq:hulog}
\end{equation}
which makes them comparable across images and suitable as classifier inputs.

\emph{Boundary descriptors.} A boundary traversed as a sequence of points yields
a chain code: each step is quantized to one of eight directions, giving a string
over an eight-symbol alphabet. The first difference of the code, taken modulo
eight, is invariant to the choice of starting point, and the minimal cyclic
rotation of the first difference is the shape number of the boundary. Because a
full shape number is a variable-length string, we use its symbol histogram,
\begin{equation}
\mathbf{c}=\Big(\frac{n_0}{N},\dots,\frac{n_7}{N}\Big),\qquad
N=\sum_{k=0}^{7}n_k,
\label{eq:chainhist}
\end{equation}
as an eight-dimensional descriptor, where $n_k$ counts occurrences of
first-difference symbol $k$. This captures the curvature mix of the boundary:
a smooth ellipse concentrates mass in a few symbols, a rough grain spreads it.

\emph{Intensity statistics.} First-order statistics of a masked region, its mean
and standard deviation, plus Shannon entropy over a 16-bin histogram,
\begin{equation}
H=-\sum_{i=1}^{16}p_i\log_2 p_i ,
\label{eq:entropy}
\end{equation}
describe how the stain distributes inside an organelle independently of its
shape, and are the natural companion to the morphological features in the cell
case.

The surveyed representations also include polygonal approximation and the
signature, which are not used here: the grain boundaries are well approximated by
an ellipse, so a polygon would add vertices without adding information, and the
signature is a plot rather than a fixed-length vector.

\subsection{Classifiers and Protocol}
\label{sec:classifiers}
Three classifier families are used, chosen to match the reference methods rather
than to maximize accuracy. Linear discriminant analysis projects onto the
direction that maximizes between-class over within-class scatter and is the
classifier of the reference cell method; Gaussian naive Bayes models each feature
independently per class and is the reference method's second comparator; the
$k$-nearest-neighbour rule and a linear support vector machine cover the grain
case, where the classes are expected to be separable by a simple boundary in
descriptor space.

All reported scores are five-fold stratified cross-validation with a fixed seed,
so every image is predicted exactly once by a model that did not see it. For the
crack case, where the measured quantity is a mask rather than a class label, the
same protocol is applied to the image-level decision rule, and the threshold
that converts a mask into a decision is selected on one half of the data and
reported on the whole subsample so that the selection is visible.

\section{Results}
\label{sec:results}

\subsection{Case 1}
\label{sec:case1}
The crack pipeline follows the thresholding-and-filtering lineage
\cite{talab2016,mohan2018}. The image is converted to luminance, smoothed, and
processed along two branches: a gradient branch computing the Sobel magnitude of
\eqref{eq:sobel}, and a threshold branch computing Otsu's threshold on the
inverted smoothed image, which makes a dark crack foreground. Each branch is
followed by opening, closing, and the removal of connected components smaller
than \crackAreaPx{} pixels, about $0.3\%$ of the $227^2$ frame.

The interesting measurement of this case is that a foreground-area decision rule
fails here, and it fails silently. If an image is declared cracked whenever its
cleaned mask covers more than a fixed fraction of the frame, every image is
declared cracked, because thresholding a concrete surface marks texture as
readily as it marks a crack: measured on the subsample, the median mask coverage
is \crackCovNeg{}\% for uncracked surfaces and \crackCovPos{}\% for cracked, so the classes
differ in coverage in the \emph{opposite} direction from the intuition that a
crack adds foreground. What separates them is not how much is marked but how it
is shaped: the median component count is \crackCompNeg{} for uncracked and
\crackCompPos{} for cracked
surfaces, and the median skeleton length of the largest component is \crackSkelNeg{}
pixels for uncracked against \crackSkelPos{} for cracked, because on a cracked surface the
crack survives cleaning as one long connected structure while on an uncracked
surface the texture survives as many short blobs.

The operating rule is therefore a classifier on the ten-element descriptor of
Section \ref{sec:features}, fitted per mask stage so that the stage ablation
compares like with like. The decision is made by a depth-four decision tree,
which is deep enough to express the shape-versus-count interaction above and
shallow enough that its splits remain readable.

Table \ref{tab:crack} ablates the three mask stages on the full subsample, with the
same depth-four rule refitted at each stage. The result is the opposite of the
usual morphological claim and is worth stating plainly: the raw threshold scores
highest, and each subsequent cleaning stage costs a fraction of a point. The
reason is that the descriptor set already contains the shape information that
cleaning is meant to supply. Skeleton length per frame diagonal carries
\crackTopImportance{} of the tree's total importance, and that quantity is
measured on the largest component, so opening, closing, and area filtering -- all
of which exist to shrink small components and fuse large ones -- change a
description that is already dominated by the largest component. On this
collection morphology is therefore clean-up that a downstream consumer would
want (a human inspector looks at a mask, not at ten numbers) but not an accuracy
gain, and reporting it as one would be a claim our measurement does not support.

\begin{table}[!t]
\centering
\caption{Case 1, mask-stage ablation on \crackImages{} images. Precision,
recall, and F1 refer to the crack class.}
\label{tab:crack}
\footnotesize
\begin{tabular}{@{}lcccc@{}}
\toprule
Mask stage & Accuracy & Precision & Recall & F1 \\
\midrule
Otsu only & 0.977 & 0.982 & 0.971 & 0.977 \\
+ opening/closing & 0.976 & 0.979 & 0.973 & 0.976 \\
+ area filter & 0.970 & 0.969 & 0.971 & 0.970 \\
\bottomrule
\end{tabular}

\end{table}

The final stage reaches accuracy \crackAcc{} with precision \crackPrec{}, recall
\crackRec{}, and F1 \crackFone{}; of the $800$ uncracked images $775$ are called uncracked and $25$ are called
cracked, and of the $800$ cracked images $23$ are called uncracked and $777$ are
called cracked. On the same subsample the Canny comparator of
Section \ref{sec:edges} reaches crack-class F1 \crackCannyFone{}, and the gradient
branch fed directly to the same descriptor-and-tree rule reaches
\crackSobelFone{}. The gradient branch thus matches the threshold branch, which
is a stronger statement about Sobel than about morphology: the $3\times3$ Sobel
kernel of \eqref{eq:sobel} already carries its own smoothing, and once the
decision is made on the largest component's skeleton rather than on total
foreground, the derivative's sensitivity to concrete texture stops mattering.
Canny is the weakest of the three because its two hysteresis thresholds are fixed
constants that were not tuned per image, which is exactly the parameter
sensitivity that motivated building the measured pipeline on Sobel in the first
place.

Crack width is measured on true positives by sampling the Euclidean distance
transform along the skeleton and doubling the mean distance, which yields
\crackWidth{} pixels. This figure is deliberately left in pixels. The collection
publishes no physical calibration, the crops carry no scale reference, and
converting to millimetres would require assuming a field of view; reporting a
millimetre value would be an invented number rather than a measurement.

The descriptor importances of the fitted tree are the direct answer to which
features are suitable here: \crackTopFeat{} carries \crackTopImportance{} of the
total importance, while roundness and eccentricity together contribute under
\crackLowImportance{}. That ordering is what the task implies. A crack is
identified by a connected structure of sufficient length, and its precise outline
is irrelevant, so descriptors that encode outline detail are dead weight while the
length and count of connected components are decisive. Logistic regression on the
same ten descriptors reaches \crackLogCV{} five-fold accuracy against
\crackTreeCV{} for the tree, so the decision boundary is nearly linear in log-area
and skeleton length, and the tree's advantage is interpretability rather than
capacity.

\subsection{Case 2}
\label{sec:case2}
The cell pipeline follows the reference method \cite{prinyakupt2015} with two
documented deviations. The reference enhances the nucleus by the colour ratio
$(R+B)/(2G)$, which is large where the nucleus is dark purple and the cytoplasm
is pale. On this collection the crops contain a black padding frame, where
$G=0$ and the ratio is undefined, so the denominator is guarded as
$(R+B+1)/(2G+1)$ and the result is clipped to its $5$th and $99$th percentiles
before scaling. The frame is additionally removed by cropping to the bounding
box of the largest bright component. The clip limits are not cosmetic: the
unguarded ratio reaches values above $50$ on a handful of specular pixels
against a median near $1$, so a $1$--$99$ clip pins both ends of the ramp and
the thresholded nucleus collapses. Histogram equalization is deliberately not
applied after the enhancement, because the ratio histogram is nearly unimodal
and equalization pushes the cytoplasm above the threshold: measured on the
subsample, the Otsu foreground share changes from $0.05$ without equalization to
$0.55$ with it, which inverts the mask.

The enhanced image is thresholded by Otsu, opened with a $3\times3$ element,
closed with a $7\times7$ element, and filtered to the largest component above
$400$ pixels, which yields the nucleus. The cell body is obtained by fitting an
ellipse to the nucleus contour and expanding it by a factor of $1.45$, followed
by a union with the nucleus so that the nucleus is a subset of the cell by
construction. This replaces the reference method's ellipse fit on a calibrated
micrograph, which this collection does not provide, and it also replaces an
intersection with a dark-pixel mask, which was measured to leave the nucleus as
almost the whole cell. Under the ellipse expansion the nucleus occupies on
average \wbcNucleusShare{} of the cell mask, and the cytoplasm is the difference
of the two.

The thirteen features of the reference method's families are computed per image:
nucleus area, perimeter, roundness, solidity, eccentricity, nucleus-to-cell area
ratio, and skeleton-length-to-area as a lobeness proxy, plus mean, standard
deviation, and 16-bin entropy of the luminance inside the nucleus and inside the
cytoplasm. Linear discriminant analysis reaches \wbcLDA{} five-fold accuracy and
naive Bayes \wbcNB{}; the better model, \wbcBest{}, is reported in full in
Table \ref{tab:wbc}, and reaches macro-F1 \wbcFone{} on \wbcImages{} images.
The confusion matrix in Table \ref{tab:wbc} is not uniform: its largest off-diagonal cell runs from eosinophils to neutrophils, followed by lymphocytes to monocytes, so the errors follow the two pairs whose nuclei are closest in size and granulation rather than spreading across all four subtypes.

\begin{table}[!t]
\centering
\caption{Case 2, per-subtype scores of the \wbcBest{} classifier under
five-fold stratified cross-validation.}
\label{tab:wbc}
\footnotesize
\begin{tabular}{@{}lcccc@{}}
\toprule
Subtype & Precision & Recall & F1 & Support \\
\midrule
Eosinophil & 0.629 & 0.360 & 0.458 & 400 \\
Lymphocyte & 0.588 & 0.542 & 0.564 & 400 \\
Monocyte & 0.613 & 0.805 & 0.696 & 400 \\
Neutrophil & 0.568 & 0.677 & 0.618 & 400 \\
\midrule
Macro average & 0.600 & 0.596 & 0.584 & 1600 \\
\bottomrule
\end{tabular}

\end{table}

The errors are structured rather than uniform, and they do not follow the
granulation axis as one might expect from the stain alone. Eosinophils and
neutrophils exchange the most images, and lymphocytes and monocytes the second
most, while monocytes are otherwise recovered well. Because the discriminating
features here are area, area ratio, and entropy rather than outline detail, a
neutrophil whose segmented nucleus merges two lobes into one component is
described almost exactly like an eosinophil, and the intensity statistics do not
separate them because both are densely granulated. This is the measured
limitation of a largest-component description on lobed nuclei.

\subsection{Case 3}
\label{sec:case3}
The grain pipeline is the classical branch of the reference work
\cite{cinar2019}: Otsu threshold on the luminance, since the grain is bright and
the background near black, closing with a $7\times7$ element to fill specular
holes, and retention of the largest connected component. An image whose second
component exceeds $60\%$ of the largest is discarded rather than resolved, since
resolving it would require a watershed step that the assignment's scope does not
include; \riceDropped{} of the sampled images were discarded this way, leaving
\riceImages{}.

Each retained grain contributes \riceFeatures{} features in three groups: nine
morphological measurements (area, perimeter, major and minor axis from the fitted
ellipse, aspect ratio, roundness of \eqref{eq:roundness}, solidity, extent,
equivalent diameter), seven Hu moments in the log form of \eqref{eq:hulog}, and
the eight-bin chain-code histogram of \eqref{eq:chainhist}.

Table \ref{tab:riceabl} is the ablation that answers the assignment's question
for this case, and the answer is negative in an informative way. Region shape
alone is the strongest group and reaches macro-F1 \riceMorphOne{} with the
better classifier; adding the seven Hu moments and the eight chain-code bins
lowers macro-F1 to \riceAllOne{}. Both added groups carry less variety
information than the nine shape measurements, so on a thousand grains they
contribute variance rather than separation: chain code alone reaches only
\riceChainOne{} and Hu moments alone \riceHuOne{}, because the first-difference histogram of a closed convex grain
is close to a constant across varieties The
reported configuration is therefore \riceBest{} on the morphological group,
with per-variety scores in Table \ref{tab:riceclasses}.

The chain code deserves a comment because it is the descriptor the course block
presents most prominently. It is the boundary descriptor par excellence when the
classes differ in outline detail, such as a dented versus a smooth grain
silhouette, and our result shows what happens when they do not: the variety
difference here is a global proportion, major over minor axis, which a
perimeter-normalized histogram of curvature symbols cannot express. The failure
is a property of the class definition, not of the descriptor.

\begin{table}[!t]
\centering
\caption{Case 3, feature-group ablation, five-fold macro-F1 on \riceImages{}
grains. Ranking is by the better of the two classifiers.}
\label{tab:riceabl}
\footnotesize
\begin{tabular*}{\columnwidth}{@{\extracolsep{\fill}}lcccc@{}}
\toprule
Feature group & kNN & LinearSVC & Best & Rank \\
\midrule
Shape (9) & 0.956 & 0.961 & 0.961 & 1 \\
Hu moments (7) & 0.803 & 0.839 & 0.839 & 3 \\
Chain code (8) & 0.457 & 0.523 & 0.523 & 4 \\
All (24) & 0.936 & 0.955 & 0.955 & 2 \\
\bottomrule
\end{tabular*}

\end{table}

Table \ref{tab:riceclasses} gives the per-variety breakdown. The residual errors
are between the two long-grain varieties, whose axis ratios overlap, while
Ipsala and Karacadag are separated almost perfectly, which again attributes the
decision to shape rather than to intensity.

\begin{table}[!t]
\centering
\caption{Case 3, per-variety scores of the \riceBest{} classifier.}
\label{tab:riceclasses}
\footnotesize
\begin{tabular}{@{}lcccc@{}}
\toprule
Variety & Precision & Recall & F1 & Support \\
\midrule
Arborio & 0.964 & 0.935 & 0.949 & 200 \\
Basmati & 0.960 & 0.950 & 0.955 & 200 \\
Ipsala & 0.995 & 1.000 & 0.998 & 200 \\
Jasmine & 0.941 & 0.955 & 0.948 & 200 \\
Karacadag & 0.946 & 0.965 & 0.955 & 200 \\
\bottomrule
\end{tabular}

\end{table}

\subsection{Computational Cost}
\label{sec:cost}
Table \ref{tab:timing} reports wall-clock time for each case on a single CPU core
with no parallelism, including feature extraction and cross-validation. The crack
case is the most expensive, and profiling attributes that cost to the per-component
descriptor stage rather than to thresholding: loading, Otsu thresholding, and the
morphological cleanup together take about $1.8$ ms per image, while the contour,
hull, and per-component distance-transform queries take about $50$ ms, because a
texture-heavy negative mask has many components. The cell case pays for the
ratio-color enhancement and three sequential morphological passes; the grain case
is cheapest because a single convex region needs one closing and one component
pass.

\begin{table}[!t]
\centering
\caption{Measured wall-clock cost per case, single core, including feature
extraction and cross-validation.}
\label{tab:timing}
\footnotesize
\begin{tabular}{@{}lccc@{}}
\toprule
Case & Images & Wall clock (s) & Per image (ms) \\
\midrule
Crack (Case 1) & 1600 & 439.3 & 274.5 \\
WBC (Case 2) & 1600 & 179.6 & 112.3 \\
Rice (Case 3) & 1000 & 3.3 & 3.3 \\
\bottomrule
\end{tabular}

\end{table}

\section{Discussion: Which Features for Which Case}
\label{sec:discussion}
Table \ref{tab:featuremap} is the paper's answer to the assignment's question. It
maps each case to its task type, the descriptor family that carried the decision
in our measurements, and the failure mode we actually observed.

\begin{table}[!t]
\centering
\caption{Feature suitability across the three cases: the descriptor family that carried the decision, and the failure mode measured in our own runs.}
\label{tab:featuremap}
\footnotesize
\begin{tabular*}{\columnwidth}{@{\extracolsep{\fill}}>{\raggedright\arraybackslash}p{0.9cm}>{\raggedright\arraybackslash}p{1.55cm}>{\raggedright\arraybackslash}p{1.9cm}>{\raggedright\arraybackslash}p{2.1cm}@{}}
\toprule
Case & Task type & Winning family & Measured failure mode \\
\midrule
Crack & detection &
skeleton length and\newline component count &
texture survives as many short components, so coverage runs opposite to intuition \\
WBC & segmentation +\newline typing &
nucleus area, area ratio, and entropy &
lobed nuclei are cut to their largest component, so a two-lobed neutrophil reads
as an eosinophil \\
Rice & classification &
ellipse axis lengths, aspect ratio, and area &
the chain-code histogram cannot express a global proportion, so boundary detail
adds variance; touching grains would force discarding rather than splitting \\
\bottomrule
\end{tabular*}
\end{table}

Three regularities hold across the cases.

First, the useful feature complexity tracks the task's decision boundary, not the
object's visual complexity. The crack decision is essentially one-dimensional, a
thin structure spanning the frame, and a depth-four tree on ten features reaches
\crackTreeCV{} five-fold accuracy, matching the hand-built cascade. The grain
decision is five-way over a global proportion, and the nine shape measurements
suffice at macro-F1 \riceMorphOne{} while the full \riceFeatures-feature vector
falls to \riceAllOne{}. Adding descriptors helped in neither case; what helped
was choosing the descriptor whose invariances match the class difference.

Second, the morphological structuring-element size is a statement about the
object's scale and must be chosen before the descriptors, not after. Opening with
a $3\times3$ element removes crack speckle without eroding the crack, but a
larger element would delete the crack itself, since the crack width measured in
Section \ref{sec:case1} is only \crackWidth{} pixels. In the cell case the
element must be large enough to bridge nuclear lobes, and the ellipse expansion
factor that defines the cell body is set by the cytoplasm's extent rather than by
the nucleus, which is why a dark-pixel intersection failed: the nucleus is the
darkest structure in the image.

Third, the descriptors that failed are informative, and one failed for a reason
worth stating. The chain-code histogram is the boundary descriptor the course
block develops in most detail, and it was the weakest grain group at
\riceChainOne{}, because a perimeter-normalized curvature histogram cannot encode
the one property that separates these varieties, the axis ratio. Harris corners
and Hough lines, surveyed in Section \ref{sec:edges}, were not used at all:
concrete cracks, cell nuclei, and rice grains are not junction- or
line-structured objects. The general lesson is that a descriptor is suitable when
the transformation group it is invariant to contains the nuisance variation of
the data and excludes the class difference, and that this can be decided from the
class definition before any classifier is trained.

\section{Conclusion}
\label{sec:conclusion}
Three classical pipelines were built for three different inspection questions and
measured on public data. Crack detection reaches accuracy \crackAcc{} and F1
\crackFone{} from gradient analysis, Otsu thresholding, morphological cleaning, and
a ten-feature shape descriptor, with a measured crack width of \crackWidth{}
pixels, and the experiment shows that the foreground-area rule which this task
usually suggests is actively wrong on real surface texture. White blood cell
subtyping reaches accuracy \wbcAcc{} with region and intensity features under a
linear discriminant, after a guarded colour-ratio enhancement of the nucleus and an
ellipse-expanded cell body. Rice variety classification reaches macro-F1
\riceMorphOne{} from nine shape measurements alone, and the ablation shows the
combined \riceFeatures-feature vector falling to \riceAllOne{}: the Hu moments and
the chain-code histogram contributed variance rather than separation because the
varieties differ in a global proportion. Across the three cases the descriptor
family that worked was predicted by the task type rather than by the object
domain, and the descriptor that failed always failed by being invariant to the
wrong transformation.

Learned feature hierarchies are out of scope here by the assignment's
constraints, and the natural next step is to use these classical descriptors as
the interpretable baseline that a convolutional model must beat, rather than
replacing them.

\begin{thebibliography}{00}

\bibitem{otsu1979} N. Otsu, ``A threshold selection method from gray-level
histograms,'' \emph{IEEE Trans. Syst., Man, Cybern.}, vol. 9, no. 1, pp. 62--66,
Jan. 1979, 10.1109/TSMC.1979.4310076.

\bibitem{canny1986} J. Canny, ``A computational approach to edge detection,''
\emph{IEEE Trans. Pattern Anal. Mach. Intell.}, vol. PAMI-8, no. 6, pp. 679--698,
Nov. 1986, 10.1109/TPAMI.1986.4767851.

\bibitem{harris1988} C. Harris and M. Stephens, ``A combined corner and edge
detector,'' in \emph{Proc. Alvey Vision Conf.}, 1988, pp. 23.1--23.6,
10.5244/C.2.23.

\bibitem{duda1972} R. O. Duda and P. E. Hart, ``Use of the Hough transformation
to detect lines and curves in pictures,'' \emph{Commun. ACM}, vol. 15, no. 1,
pp. 11--15, Jan. 1972, 10.1145/361237.361242.

\bibitem{talab2016} A. M. A. Talab, Z. Huang, F. Xi, and L. HaiMing,
``Detection crack in image using Otsu method and multiple filtering in image
processing techniques,'' \emph{Optik}, vol. 127, no. 3, pp. 1030--1033, Feb.
2016, 10.1016/j.ijleo.2015.09.147.

\bibitem{prinyakupt2015} J. Prinyakupt and C. Pluempitiwiriyawej,
``Segmentation of white blood cells and comparison of cell morphology by linear
and na\"ive Bayes classifiers,'' \emph{BioMed. Eng. OnLine}, vol. 14, art. 63,
Jun. 2015, 10.1186/s12938-015-0037-1.

\bibitem{cinar2019} I. Cinar and M. Koklu, ``Classification of rice varieties
using artificial intelligence methods,'' \emph{Int. J. Intell. Syst. Appl. Eng.},
vol. 7, no. 3, pp. 188--194, Sep. 2019, 10.18201/ijisae.2019355381.

\bibitem{mohan2018} A. Mohan and S. Poobal, ``Crack detection using image
processing: A critical review and analysis,'' \emph{Alexandria Eng. J.}, vol. 57,
no. 2, pp. 787--798, Jun. 2018, 10.1016/j.aej.2017.01.020.

\bibitem{ozgenel2019} \c{C}. F. \"Ozgenel, ``Concrete crack images for
classification,'' Mendeley Data, V2, 2019, 10.17632/5y9wdsg2zt.2.

\bibitem{dorafshan2018} S. Dorafshan, R. J. Thomas, and M. Maguire,
``SDNET2018: An annotated image dataset for non-contact concrete crack detection
using deep convolutional neural networks,'' \emph{Data in Brief}, vol. 21,
pp. 1664--1668, Dec. 2018, 10.1016/j.dib.2018.11.015.

\bibitem{koklu2021} M. Koklu, I. Cinar, and Y. S. Taspinar, ``Classification of
rice varieties with deep learning methods,'' \emph{Comput. Electron. Agric.},
vol. 187, art. 106285, Aug. 2021, 10.1016/j.compag.2021.106285.

\bibitem{mooney2018} P. Mooney, ``Blood cell images,'' Kaggle dataset, accessed
Nov. 13, 2026. [Online]. Available:
https://www.kaggle.com/paultimothymooney/blood-cells

\bibitem{gonzalez2018} R. C. Gonzalez and R. E. Woods, \emph{Digital Image
Processing}, 4th ed. New York, NY, USA: Pearson, 2018.

\bibitem{handouts} Wahyono, ``Morphological processing,'' ``Image feature
extraction (edge detection),'' ``Image feature extraction (corner and line),''
and ``Feature representation (edge-based),'' Digital Image Processing course
handouts, Dept. Comput. Sci. Electron., Universitas Gadjah Mada, Yogyakarta,
Indonesia, n.d.

\end{thebibliography}

\EOD

\end{document}
